% =====================================================================
%  tikz.tex  --  every drawing in the deck, as a macro.
%  Loaded from main.tex BEFORE \begin{document}, so no parameterised
%  tikz style is ever defined inside a frame.
%
%  Chart data: figures quoted in the text of Lester et al. (2021) are
%  exact (71.8, 89.3, and every parameter count, from the prose and
%  Table 4).  Curve values that exist in that paper only as plotted
%  lines are digitised approximations -- good enough to carry the
%  argument, not citable to a decimal place.
% =====================================================================

% ---------------------------------------------------------------------
% \ptsnow{(x,y)} : tiny frost mark meaning "frozen / not updated"
% ---------------------------------------------------------------------
\newcommand{\ptsnow}[1]{%
  \begin{scope}[shift={#1}]
    \foreach \a in {0,60,120}{\draw[line width=0.6pt,draw=ptblue] (\a:0.13) -- (\a+180:0.13);}
  \end{scope}}

% ---------------------------------------------------------------------
% \ptdot{color} : small inline colour chip, usable inside running text
% ---------------------------------------------------------------------
\newcommand{\ptdot}[1]{%
  \tikz[baseline=-0.55ex]\node[ptchip,draw=#1,fill=#1!18,
        minimum width=0.42cm,minimum height=0.30cm]{};}

% ---------------------------------------------------------------------
% \ptprompttuning : frozen model fed by a learned soft prompt (chips)
% plus the ordinary input text (frozen embeddings) -- half-width column.
% Coordinate map: model box (0,2.55) 4.3x1.15;
% soft-prompt chips centred x=-1.15, 3 chips, 0.5cm pitch, y=0;
% input chips centred x=1.75, 6 chips, 0.44cm pitch, y=0;
% braces y=-0.26, labels y=-0.65.
% ---------------------------------------------------------------------
\newcommand{\ptprompttuning}{%
  \begin{tikzpicture}

    \node[ptfrozen,minimum width=4.3cm,minimum height=1.15cm,align=center,font=\small]
          (lm) at (0,2.55) {Pre-trained Model\\[1pt]{\bfseries\footnotesize Frozen}};
    \ptsnow{(1.75,2.15)}

    \foreach \i in {0,1,2}{
      \node[fill=ptgreen!18,draw=ptgreen,line width=0.9pt,rounded corners=3pt,
            minimum width=0.46cm,minimum height=0.46cm] at (-1.55+0.5*\i,0) {};}
    \draw[ptarrow] (-1.15,0.28) -- (-1.15,1.92);

    \foreach \i in {0,...,5}{
      \node[fill=white,draw=ptrule,line width=0.8pt,rounded corners=3pt,
            minimum width=0.44cm,minimum height=0.44cm] at (1.10+0.44*\i,0) {};}
    \draw[ptarrow] (1.75,0.28) -- (1.75,1.92);

    \draw[decorate,decoration={brace,amplitude=4pt,mirror},draw=ptrule]
         (-1.75,-0.26) -- (-0.55,-0.26);
    \node[ptlbl,anchor=north,align=center,font=\scriptsize,text=ptgreen]
         at (-1.15,-0.68) {Tunable\\Soft Prompt};

    \draw[decorate,decoration={brace,amplitude=4pt,mirror},draw=ptrule]
         (0.88,-0.26) -- (2.62,-0.26);
    \node[ptlbl,anchor=north,font=\scriptsize,text=ptmuted]
         at (1.75,-0.68) {Input Text};

  \end{tikzpicture}}


% ---------------------------------------------------------------------
% \ptanalogy : written instruction vs its learned counterpart
% ---------------------------------------------------------------------
\newcommand{\ptanalogy}{%
  \begin{tikzpicture}

    % ---------------------------------------------------------------
    % Overlay 1: only the left box + its label, centered
    % ---------------------------------------------------------------
    \only<1>{%
      \node[
        ptnode,
        minimum width=4.4cm,
        minimum height=1.9cm,
        text width=3.9cm,
        font=\small,
        align=left
      ] (a1) at (5.25,0)
      {``Write a poem with an ABAB rhyme scheme, in iambic tetrameter, about love.''};

      \node[
        ptlbl,
        anchor=south
      ] at (5.25,1.05)
      {what we would write};
    }

    % ---------------------------------------------------------------
    % Overlay 2: original complete layout
    % ---------------------------------------------------------------
    \only<2>{%
      \node[
        ptnode,
        minimum width=4.4cm,
        minimum height=1.9cm,
        text width=3.9cm,
        font=\small,
        align=left
      ] (a) at (2.30,0)
      {``Write a poem with an ABAB rhyme scheme, in iambic tetrameter, about love.''};

      \node[
        pttuned,
        minimum width=4.4cm,
        minimum height=1.9cm,
        text width=3.9cm,
        align=center
      ] (b) at (8.20,0)
      {\ttfamily\scriptsize
        0.14\ \ $-$0.72\ \ 0.31\ \ 0.08\\[3pt]
        $-$0.45\ \ 0.19\ \ 0.66\ \ $-$0.27\\[3pt]
        0.52\ \ $-$0.11\ \ 0.03\ \ 0.94
      };

      \node[ptlbl,anchor=south] at (2.30,1.05)
        {what we would write};

      \node[ptlblg,anchor=south] at (8.20,1.05)
        {what it learns instead};

      \draw[ptarrow] (4.60,0) -- (5.90,0);
    }

  \end{tikzpicture}%
}

% ---------------------------------------------------------------------
% \ptbenefitscycle : four-arc pinwheel + 2x2 label grid (Benefits frame)
% Coordinate map: ring r=1.5 centred at ORIGIN, shifted left by \ringshift.
% Left column right-edges anchor at x=-2.2, right column left-edges
% anchor at x=2.2 -- perfectly symmetric around the (unshifted) centre.
% Each title is wrapped in \mbox{} so "chip + Title" can never hyphenate
% mid-word; text width is generous enough that it always fits on one line.
% Arcs use a thinner stroke (5pt) with an oversized Stealth tip so the
% arrowhead reads clearly instead of being absorbed by a round cap.
% ---------------------------------------------------------------------
\newcommand{\ptbenefitscycle}{%
  \begin{tikzpicture}
    \def\R{1.5}
    \def\ringshift{-0.4}  % negative = left; tweak this number to taste

    \begin{scope}[shift={(\ringshift,0)}]
      \foreach \c/\col in {45/ptgreen,135/ptorange,225/ptgold,315/ptblue}{
        \draw[-{Stealth[length=9pt,width=8pt]},line width=5pt,line cap=round,draw=\col]
             ({\c+32}:\R) arc ({\c+32}:{\c-32}:\R);}
    \end{scope}

    % top-left : Scale Solutions (orange)
    \node[anchor=east,align=left,text width=4.6cm] at (-2.2,1.55)
      {\mbox{\ptdot{ptorange}\ \textbf{\color{ptorange}Scale Solutions}}\\[3pt]
       {\footnotesize\color{ptmuted}Implement AI across industries efficiently.}};

    % bottom-left : Reduce Costs (gold)
    \node[anchor=east,align=left,text width=4.6cm] at (-2.2,-1.55)
      {\mbox{\ptdot{ptgold}\ \textbf{\color{ptgold}Reduce Costs}}\\[3pt]
       {\footnotesize\color{ptmuted}Lower resource usage and expenses.}};

    % top-right : Optimize Prompts (green)
    \node[anchor=west,align=left,text width=4.6cm] at (2.2,1.55)
      {\mbox{\ptdot{ptgreen}\ \textbf{\color{ptgreen}Optimize Prompts}}\\[3pt]
       {\footnotesize\color{ptmuted}Refine input prompts for better results.}};

    % bottom-right : Enhance Accuracy (blue)
    \node[anchor=west,align=left,text width=4.6cm] at (2.2,-1.55)
      {\mbox{\ptdot{ptblue}\ \textbf{\color{ptblue}Enhance Accuracy}}\\[3pt]
       {\footnotesize\color{ptmuted}Improve model accuracy and context awareness.}};
  \end{tikzpicture}}

% ---------------------------------------------------------------------
% \ptfullfinetune : Original LLM --(update all weights)--> Full Fine-Tuned LLM
% ---------------------------------------------------------------------
\newcommand{\ptfullfinetune}{%
  \begin{tikzpicture}
    \node[circle,draw=ptorange,fill=ptorange!15,line width=1.1pt,
          minimum size=2.3cm,align=center,font=\small\bfseries,text=ptink]
          (orig) at (0,0) {Original\\LLM};
    \node[circle,draw=ptorange,fill=ptorange!15,line width=1.1pt,
          minimum size=2.3cm,align=center,font=\scriptsize\bfseries,text=ptink]
          (ft) at (6.4,0) {Full\\Fine-Tuned\\LLM};
    \draw[ptarrow] (orig.east) -- (ft.west)
          node[midway,above=2pt,align=center,font=\scriptsize,text=ptink]
          {Update all the\\weights during\\training};
  \end{tikzpicture}}

% ---------------------------------------------------------------------
% \pttuningcompare : model tuning (left) vs prompt tuning (right)
% ---------------------------------------------------------------------
\newcommand{\pttuningcompare}{%
  \begin{tikzpicture}
    \node[ptlbl,font=\small\bfseries,text=ptblue]  at (1.85,2.20) {Model tuning};
    \node[ptlbl,font=\small\bfseries,text=ptgreen] at (7.40,2.20) {Prompt tuning};
    \foreach \n/\y in {A/1.15,B/0,C/-1.15}{
      \node[ptghost,minimum width=1.3cm,minimum height=0.62cm,font=\scriptsize] (l\n) at (0.45,\y) {Task \n};
      \node[ptnode,fill=ptblue!10,draw=ptblue,minimum width=2.5cm,minimum height=0.72cm,font=\scriptsize]
            (m\n) at (3.05,\y) {Task \n\ model, 11B};
      \draw[ptarrow] (l\n.east) -- (m\n.west);}
    \draw[ptrule,dashed,line width=0.9pt] (4.85,-2.30) -- (4.85,2.20);
    \node[ptghost,minimum width=1.5cm,minimum height=2.3cm] (mix) at (6.00,0) {};
    \foreach \i/\c in {0/ptgreen,1/ptorange,2/ptgold}{
      \node[ptchip,draw=\c,fill=\c!18,minimum width=0.5cm,minimum height=0.34cm]
            at (6.00,0.62-0.62*\i) {};}
    \node[ptfrozen,minimum width=2.5cm,minimum height=2.3cm,align=center] (fm) at (8.75,0)
         {frozen T5\\[-2pt]{\scriptsize 11B, shared}};
    \ptsnow{(8.75,-0.78)}
    \draw[ptarrow] (mix.east) -- (fm.west);
    \node[ptlbl,anchor=north] at (6.00,-1.35) {mixed batch};
    \node[ptlbl,anchor=north,text width=4.6cm] at (7.40,-2.00)
         {one model, one batch, three tasks at once};
    \node[ptlbl,anchor=north,text width=4.4cm] at (1.85,-1.65)
         {a copy per task, a batch per task};
  \end{tikzpicture}}

% ---------------------------------------------------------------------
% \ptinitmanifold : why TEXT/SAMPLE_VOCAB land inside the manifold and
% RANDOM can land outside it. An informal blob, not a real geometric
% embedding -- it only needs to carry "inside" vs "outside".
% ---------------------------------------------------------------------
\newcommand{\ptinitmanifold}{%
  \begin{tikzpicture}
    \fill[ptpanel,rounded corners=40pt] (-3.4,-1.7) rectangle (3.4,1.7);
    \draw[ptrule,rounded corners=40pt,line width=0.8pt] (-3.4,-1.7) rectangle (3.4,1.7);
    \node[ptlbl,anchor=north,font=\scriptsize] at (0,-2.05)
         {the embedding manifold --- where real token vectors live};
    \fill[ptgreen] (-1.6,0.5) circle (3.2pt);
    \node[anchor=west,font=\scriptsize\bfseries,text=ptgreen] at (-1.35,0.5) {TEXT};
    \fill[ptgold] (1.1,-0.55) circle (3.2pt);
    \node[anchor=west,font=\scriptsize\bfseries,text=ptgold] at (1.35,-0.55) {SAMPLE\_VOCAB};
    \draw[ptrule,dashed,line width=0.8pt] (3.4,1.15) -- (4.45,1.15);
    \fill[ptred] (4.6,1.15) circle (3.2pt);
    \node[anchor=west,font=\scriptsize\bfseries,text=ptred] at (4.85,1.15) {RANDOM};
  \end{tikzpicture}}

% ---------------------------------------------------------------------
% \ptconverge : loss vs training steps for the soft prompt, with a
% dashed reference line for where full model tuning typically lands.
% Values are illustrative, matching the paper's reported 30,000-step
% training schedule (Appendix A.1) -- not digitised from any figure.
% ---------------------------------------------------------------------
\newcommand{\ptconverge}{%
\begin{tikzpicture}
\begin{axis}[ptaxis,xlabel={training steps},ylabel={loss},
             width=9.2cm,height=5.0cm,xmin=0,xmax=30000,ymin=0,ymax=3.6,
             xtick={0,10000,20000,30000},xticklabels={0,10K,20K,30K},
             legend pos=north east,legend cell align=left]
  \addplot[ptcurve,color=ptgreen,mark=none,line width=1.6pt] coordinates
    {(0,3.4)(2000,2.6)(5000,1.9)(10000,1.3)(15000,0.95)(20000,0.75)(25000,0.62)(30000,0.55)};
  \addlegendentry{soft prompt loss}
  \draw[ptrule,dashed,line width=0.9pt] (axis cs:0,0.48) -- (axis cs:30000,0.48);
  \node[anchor=north east,font=\scriptsize,text=ptmuted] at (axis cs:29500,0.62)
       {model tuning floor};
\end{axis}
\end{tikzpicture}}

% ---------------------------------------------------------------------
% \pttrainloop : the training cycle
% Coordinate map:
%   (init) (1.30, 1.00)  2.6 x 0.9
%   (fwd)  (4.80, 1.00)  2.8 x 0.9
%   (loss) (8.30, 1.00)  2.6 x 0.9
%   (upd)  (6.55,-0.95)  3.6 x 0.9
%   return: loss.south -> (8.30,-0.95) -> upd.east ; upd.west -> (4.80,-0.95) -> fwd.south
% ---------------------------------------------------------------------
\newcommand{\pttrainloop}{%
  \begin{tikzpicture}
    % Top row (visible on both slides)
    \node[pttuned,minimum width=2.6cm,minimum height=0.9cm,font=\scriptsize] (init) at (1.30,1.00)
         {initialise the\\[-1pt]soft prompt};

    \node[ptnode,minimum width=2.8cm,minimum height=0.9cm,font=\scriptsize] (fwd) at (4.80,1.00)
         {forward pass\\[-1pt]through the frozen model};

    \node[ptnode,minimum width=2.6cm,minimum height=0.9cm,font=\scriptsize] (loss) at (8.30,1.00)
         {cross-entropy\\[-1pt]loss};

    % Top arrows (visible on both slides)
    \draw[ptarrow] (init.east) -- (fwd.west);
    \draw[ptarrow] (fwd.east)  -- (loss.west);

    % Bottom training loop (appears only on overlay 2)
    \only<2->{%
        \node[pttuned,minimum width=3.6cm,minimum height=1.1cm,font=\scriptsize] 
        (upd) at (6.55,-0.95)
        {backpropagate,\\[0.0pt]update the soft prompt only};

        \draw[ptarrow] (loss.south) -- (loss.south |- upd.north);
        \draw[ptarrow] (fwd.south |- upd.north) -- (fwd.south);

        \node[ptlbl,anchor=north] at (3.00,-0.7) 
        {repeat until convergence};
    }
  \end{tikzpicture}%
}

% ---------------------------------------------------------------------
% \ptsoftpromptinit : four learnable prompt-vector chips, prepended to
% an (untouched) input-text block, with a curly brace under the row
% instead of an arrow pointing into the middle of it.
% ---------------------------------------------------------------------
\newcommand{\ptsoftpromptinit}{%
  \begin{tikzpicture}
    \foreach \i/\lbl in {0/p_1,1/p_2,2/p_3,3/p_4}{
      \node[pttuned,minimum width=0.85cm,minimum height=0.85cm,font=\small]
            (p\i) at (0.95*\i,0) {$\lbl$};}
    \node[font=\large,text=ptmuted] at (3.55,0) {$+$};
    \node[ptghost,minimum width=2.6cm,minimum height=0.85cm,font=\small]
          (inp) at (5.35,0) {Input Text};
    \draw[decorate,decoration={brace,amplitude=6pt,mirror},draw=ptgreen,line width=0.9pt]
         (-0.425,-0.55) -- (3.275,-0.55);
    \node[ptlblg,anchor=north] at (1.425,-0.95) {Learnable};
  \end{tikzpicture}}

% ---------------------------------------------------------------------
% \ptspectrum : Prompt Engineering -> Prompt Tuning -> Fine-Tuning
% ---------------------------------------------------------------------
\newcommand{\ptspectrum}{%
  \begin{tikzpicture}
    \draw[ptrule, line width=1.4pt, -{Stealth[length=6pt]}] (0,0) -- (10.4,0);
    \foreach \x/\lbl/\col/\sub in {
        0.4/{Prompt\\Engineering}/ptslate/{write it by hand},
        5.2/{Prompt\\Tuning}/ptgreen/{learn a few vectors},
        10.0/{Fine-Tuning}/ptorange/{update every weight}}{
      \fill[\col] (\x,0) circle (3.2pt);
      \node[ptlbl,font=\small\bfseries,text=\col,align=center] at (\x,0.75) {\lbl};
      \node[ptlbl,font=\scriptsize,align=center] at (\x,-0.55) {\sub};
    }
    \node[ptlbl,font=\scriptsize,text=ptmuted] at (5.2,-1.35)
         {\(\longleftarrow\) faster, cheaper \hspace{1.4cm} more control, more cost \(\longrightarrow\)};
  \end{tikzpicture}}

% ---------------------------------------------------------------------
% \ptradar : Fine-Tuning vs Prompt Tuning, 5-axis radar chart (diagram
% only -- the legend is built separately in the slide, to its left).
% ---------------------------------------------------------------------
\newcommand{\ptradar}{  \begin{tikzpicture}    \def\R{1.75cm}     grid rings    \foreach \r in {1,2,3,4}{      \draw[ptrule,dashed,line width=0.3pt] (0,0) circle ({\r*\R/4});}    \draw[ptrule,line width=0.5pt] (0,0) circle (\R);     axes    \foreach \ang in {90,18,-54,-126,162}{      \draw[ptrule,line width=0.5pt] (0,0) -- (\ang:\R);}     labels: fixed offset beyond the ring, manual line breaks, anchor    chosen per side so the text grows outward, never back over the ring    \node[ptlbl,font=\small,align=center,anchor=south]         at (90:{\R+0.75cm}) {Speed of\\Implementation};    \node[ptlbl,font=\small,align=left,anchor=west]         at (18:{\R+0.60cm}) {Data\\Requirement};    \node[ptlbl,font=\small,align=left,anchor=north west]         at (-54:{\R+0.55cm}) {Output\\Quality};    \node[ptlbl,font=\small,align=right,anchor=north east]         at (-126:{\R+0.55cm}) {Model\\Modification};    \node[ptlbl,font=\small,align=right,anchor=east]         at (162:{\R+0.60cm}) {Flexibility};     Fine-tuning polygon: Speed 3, Data 9, Output 8, ModMod 8, Flex 5    \draw[ptorange,line width=1.3pt,fill=ptorange,fill opacity=0.14]      (90:{3/9*\R}) -- (18:{9/9*\R}) -- (-54:{8/9*\R}) --      (-126:{8/9*\R}) -- (162:{5/9*\R}) -- cycle;     Prompt tuning polygon: Speed 8, Data 2, Output 4, ModMod 1, Flex 7    \draw[ptslate,line width=1.3pt,fill=ptslate,fill opacity=0.18]      (90:{8/9*\R}) -- (18:{2/9*\R}) -- (-54:{4/9*\R}) --      (-126:{1/9*\R}) -- (162:{7/9*\R}) -- cycle;  \end{tikzpicture}}

% ---------------------------------------------------------------------
% \ptwexinit : worked-example Step 2 -- concatenate p1,p2,p3 with the
% first training example, then collapse down into "the soft prompt".
% Coordinate map: chips at x=0,0.95,1.9, y=0; plus sign x=2.55;
% input box centred x=5.55, y=0; arrow down from x=2.75 to y=-1.3;
% result label at y=-1.7.
% ---------------------------------------------------------------------
% ---------------------------------------------------------------------
% \ptwexinit : worked-example Step 2 -- concatenate p1,p2,p3 with the
% first training example, then collapse down into "the soft prompt".
% Coordinate map: chips at x=0,0.95,1.9, y=0; plus sign x=2.55;
% input box centred x=5.55, y=0; arrow down from x=2.75 to y=-1.3;
% result label at y=-2.65.
% ---------------------------------------------------------------------
% ---------------------------------------------------------------------
% \ptwexinit : worked-example Step 2 -- concatenate p1,p2,p3 with the
% first training example, then collapse down into "the soft prompt".
% Coordinate map: chips at x=0,0.95,1.9, y=0; plus sign x=2.55;
% input box centred x=5.55, y=0; brace at y=-0.55; arrow y=-0.75 to -1.30;
% result box at y=-1.75.
% ---------------------------------------------------------------------
\newcommand{\ptwexinit}{%
  \begin{tikzpicture}

    % ==================================================
    % Overlay 1: Only the input sentence, solid style
    % Same size as the input box on Overlay 2
    % ==================================================
    \only<1>{%
      \node[
        ptnode,
        minimum width=4.3cm,
        minimum height=0.85cm,
        font=\small
      ] at (3.55,0)
      {``The movie was fantastic.''};
    }

    % ==================================================
    % Overlay 2 onwards: Full soft prompt illustration
    % ==================================================
    \only<2->{%

      % Soft prompt tokens
      \foreach \i/\lbl in {0/p_1,1/p_2,2/p_3}{
        \node[
          pttuned,
          minimum width=0.85cm,
          minimum height=0.85cm,
          font=\small
        ]
        (p\i) at (0.95*\i,0) {$\lbl$};
      }

      % Plus sign
      \node[
        font=\large,
        text=ptmuted
      ] (plus) at (2.89,0) {$+$};

      % Original input sentence
      \node[
        ptghost,
        minimum width=4.3cm,
        minimum height=0.85cm,
        font=\small
      ]
      (inp) at (5.55,0)
      {``The movie was fantastic.''};

      % Brace under prompt tokens
      \draw[
        decorate,
        decoration={brace,amplitude=6pt,mirror},
        draw=ptgreen,
        line width=0.9pt
      ]
      (-0.425,-0.55) -- (2.325,-0.55);

      % Arrow to label
      \draw[
        ptarrow,
        line width=1.0pt
      ]
      (0.95,-0.85) -- (0.95,-1.35);

      % Soft Prompt label
      \node[
        pttuned,
        minimum width=2.6cm,
        minimum height=0.68cm,
        font=\small\bfseries
      ]
      (sp) at (0.95,-1.75)
      {Soft Prompt};
    }

  \end{tikzpicture}%
}

  % ---------------------------------------------------------------------
% \ptwexfreeze : worked-example Step 3 -- frozen pretrained model above,
% soft prompt P below, connected by a "weights frozen" arrow, with a
% small checkmark note that only P is updated.
% Coordinate map: model box (0,1.4) 4.6x1.4; snow mark at (2.0,1.05);
% arrow x=0, y=0.55 down to y=-0.35; label at y=0.10;
% soft prompt box (0,-0.9) 3.4x0.85; "updated" label at y=-1.45.
% ---------------------------------------------------------------------
% ---------------------------------------------------------------------
% \ptwexfreeze : worked-example Step 3 -- animated version
%
% Overlay 1:
%   - Shows only the "Pretrained Model Frozen" box
%   - Snow mark remains visible
%   - Everything below the model is hidden
%
% Overlay 2:
%   - Shows the complete diagram
%   - Arrow, "weights frozen" label, soft prompt, and
%     "only P is updated" note appear
% ---------------------------------------------------------------------
% ---------------------------------------------------------------------
% \ptwexfreeze : worked-example Step 3 -- animated version
% ---------------------------------------------------------------------
\newcommand{\ptwexfreeze}{%
  \begin{tikzpicture}

    % Pretrained model -- visible on both overlays
    \node[
      ptfrozen,
      minimum width=4.6cm,
      minimum height=1.4cm,
      align=center,
      font=\small
    ]
    (lm) at (0,1.4)
    {Pretrained Model\\[1pt]{\bfseries\footnotesize Frozen}};

    \ptsnow{(1.95,1.05)}

    % ================================================================
    % Bottom section -- hidden on overlay 1, visible on overlay 2
    % Space is preserved on both overlays
    % ================================================================
    \uncover<2->{%

      % Arrow showing frozen weights
      \draw[
        ptarrow,
        line width=1.0pt
      ]
      (0,0.55) -- (0,-0.30);

      % "weights frozen" label
      \node[
        ptlbl,
        anchor=west,
        font=\scriptsize,
        text=ptmuted
      ]
      at (0.25,0.13)
      {weights frozen};

      % Soft prompt
      \node[
        pttuned,
        minimum width=3.4cm,
        minimum height=0.85cm,
        font=\small\bfseries
      ]
      (sp) at (0,-0.75)
      {$P=[p_1,p_2,p_3]$};

      % Updated note
      \node[
        ptlblg,
        anchor=north,
        font=\scriptsize
      ]
      at (0,-1.20)
      {only $P$ is updated};

    }

  \end{tikzpicture}%
}


  % ---------------------------------------------------------------------
% \ptwexforward : worked-example Step 4 -- [p1 p2 p3] + input flows into
% the frozen model, producing a two-class probability readout where the
% predicted class (Negative, 0.60) is wrong against the true label
% (Positive, 0.40), flagged in ptred.
% Coordinate map: chip/input row y=2.0; model box (0,0.9) 4.6x1.0;
% arrows at x=0; prediction row y=-0.55, two boxes at x=-1.3 and x=1.3.
% ---------------------------------------------------------------------
% ---------------------------------------------------------------------
% \ptwexforward : worked-example Step 4 -- animated version
% ---------------------------------------------------------------------
\newcommand{\ptwexforward}{%
  \begin{tikzpicture}

    % ================================================================
    % Top row -- kept exactly as before
    % ================================================================
    \foreach \i/\lbl in {0/p_1,1/p_2,2/p_3}{
      \node[
        pttuned,
        minimum width=0.72cm,
        minimum height=0.72cm,
        font=\small
      ]
      (p\i) at (-2.05+0.78*\i,2.5) {$\lbl$};
    }

    \node[
      ptghost,
      minimum width=4.0cm,
      minimum height=0.72cm,
      font=\small
    ]
    (inp) at (2.3,2.5)
    {``The movie was fantastic.''};

    % ================================================================
    % Central vertical arrow
    % Aligned below the top row
    % ================================================================
    \draw[
      ptarrow,
      line width=1.0pt
    ]
    (0.8,1.82) -- (0.8,1.42);

    % ================================================================
    % Frozen model
    % Centered underneath the top row
    % ================================================================
    \node[
      ptfrozen,
      minimum width=4.6cm,
      minimum height=1.0cm,
      align=center,
      font=\small
    ]
    (lm) at (0.8,0.85)
    {Frozen Model};

    % Snowflake / frozen indicator
    \ptsnow{(2.65,0.85)}

    % ================================================================
    % Overlay 2: prediction section
    % Animation strategy unchanged
    % ================================================================
    \uncover<2->{%

      % --------------------------------------------------------------
      % Central downward arrow
      % Aligned with Frozen Model
      % --------------------------------------------------------------
      \draw[
        ptarrow,
        line width=1.0pt
      ]
      (0.8,0.15) -- (0.8,-0.25);

      % --------------------------------------------------------------
      % Prediction label
      % --------------------------------------------------------------
      \node[
        ptlbl,
        anchor=west,
        font=\scriptsize,
        text=ptmuted
      ]
      at (1.00,-0.05)
      {prediction};

      % --------------------------------------------------------------
      % Probability boxes
      % Centered around x=0.8
      % --------------------------------------------------------------
      \node[
        ptnode,
        fill=white,
        draw=ptrule,
        minimum width=2.2cm,
        minimum height=0.68cm,
        font=\scriptsize
      ]
      (pos) at (-0.50,-0.75)
      {Positive: $0.40$};

      \node[
        ptnode,
        fill=ptred!12,
        draw=ptred,
        minimum width=2.2cm,
        minimum height=0.68cm,
        font=\scriptsize\bfseries,
        text=ptred
      ]
      (neg) at (2.10,-0.75)
      {Negative: $0.60$};

      % --------------------------------------------------------------
      % Actual class
      % --------------------------------------------------------------
      \node[
        ptlbl,
        anchor=north,
        font=\scriptsize,
        text=ptmuted
      ]
      at (-0.50,-1.5)
      {actual class};

      % --------------------------------------------------------------
      % Predicted class
      % --------------------------------------------------------------
      \node[
        ptlbl,
        anchor=north,
        font=\scriptsize,
        text=ptmuted
      ]
      at (2.10,-1.5)
      {predicted class};

    }

  \end{tikzpicture}%
}

  % ---------------------------------------------------------------------
% \ptwexloss : worked-example Step 5 -- take the Positive probability
% (0.40) from Step 4, run it through the cross-entropy formula, and
% land on the numeric loss value.
% Coordinate map: prob box (0,1.3); formula row y=0.35;
% substitution row y=-0.35; result stat at y=-1.15.
% ---------------------------------------------------------------------
% ---------------------------------------------------------------------
% \ptwexloss : worked-example Step 5 -- take the Positive probability
% (0.40) from Step 4, run it through the cross-entropy formula, and
% land on the numeric loss value.
% ---------------------------------------------------------------------
% ---------------------------------------------------------------------
% \ptwexloss : worked-example Step 5 -- take the Positive probability
% (0.40) from Step 4, run it through the cross-entropy formula, and
% land on the numeric loss value.
% ---------------------------------------------------------------------
\newcommand{\ptwexloss}{%
  \begin{tikzpicture}

    % ---------------------------------------------------------------
    % Positive probability
    % ---------------------------------------------------------------
    \node[
      ptnode,
      fill=white,
      draw=ptrule,
      minimum width=3.6cm,
      minimum height=0.75cm,
      font=\small
    ]
    (prob) at (0,0)
    {Positive probability: $0.40$};

    % ---------------------------------------------------------------
    % Arrow to formula
    % ---------------------------------------------------------------
    \draw[
      ptarrow,
      line width=1.0pt
    ]
    (0,-0.4) -- (0,-0.7);

    % ---------------------------------------------------------------
    % Cross-entropy formula
    % ---------------------------------------------------------------
    \node[
      font=\normalsize,
      text=ptink
    ]
    at (0,-1.0)
    {$L = -\log(p_{\text{correct}})$};

    % ---------------------------------------------------------------
    % Substitution
    % ---------------------------------------------------------------
    \node[
      font=\normalsize,
      text=ptink
    ]
    at (0,-1.65)
    {$L = -\log(0.40)$};

    % ---------------------------------------------------------------
    % Overlay 2:
    % Arrow and final numerical loss box appear only on slide 2.
    % \uncover keeps their space/position unchanged on slide 1.
    % ---------------------------------------------------------------
    \uncover<2->{%

      \draw[
        ptarrow,
        line width=1.0pt
      ]
      (0,-1.95) -- (0,-2.25);

      \node[
        pttuned,
        minimum width=3cm,
        minimum height=0.8cm,
        font=\large\bfseries
      ]
      (loss) at (0,-2.75)
      {$L \approx 0.916$};

    }

  \end{tikzpicture}%
}

  % ---------------------------------------------------------------------
% \ptwexbackprop : worked-example Step 6 -- loss flows backward through
% the frozen model, gradients reach the soft-prompt chips, and the
% optimizer picks up nabla_P L. Frozen model uses ptsnow, chips use
% pttuned, matching Steps 2/3/4.
% Coordinate map: loss box y=1.7; frozen model y=0.45; chips y=-0.85;
% grad label y=-1.25; optimizer ghost box y=-2.15.
% ---------------------------------------------------------------------
% ---------------------------------------------------------------------
% \ptwexbackprop : worked-example Step 6 -- loss flows backward through
% the frozen model, gradients reach the soft-prompt chips, and the
% optimizer picks up nabla_P L. Compact vertical spacing so it fits
% next to a text column without overflowing the frame.
% Coordinate map: loss box y=1.15; frozen model y=0.30; chips y=-0.60;
% grad label y=-0.92; optimizer ghost box y=-1.55.
% ---------------------------------------------------------------------
% ---------------------------------------------------------------------
% \ptwexbackprop : worked-example Step 6 -- loss flows backward through
% the frozen model, gradients reach the soft-prompt chips, and the
% optimizer picks up nabla_P L. Every gap below accounts for the actual
% node half-heights, so nothing overlaps once resized.
% Coordinate map: loss y=2.0 (h0.55); model y=0.80 (h0.85); chips
% y=-0.365 (h0.58); grad label y=-0.805; optimizer y=-1.43 (h0.45).
% ---------------------------------------------------------------------
\newcommand{\ptwexbackprop}{%
  \begin{tikzpicture}
    \node[ptnode,fill=white,draw=ptrule,minimum width=1.8cm,minimum height=0.55cm,font=\small\bfseries]
          (loss) at (0,2.0) {Loss};

    \draw[ptarrow,line width=1.0pt] (0,1.62) -- (0,1.32);
    \node[ptlbl,anchor=west,font=\scriptsize,text=ptmuted] at (0.18,1.50);

    \node[ptfrozen,minimum width=3.0cm,minimum height=0.85cm,align=center,font=\small]
          (lm) at (0,0.80) {Frozen Model};
    \ptsnow{(1.65,0.80)}

    \draw[ptarrow,line width=1.0pt] (0,0.3) -- (0,-0.02);
    \node[ptlbl,anchor=west,font=\scriptsize,text=ptmuted] at (0.18,0.13) ;

    \foreach \i/\lbl in {0/p_1,1/p_2,2/p_3}{
      \node[pttuned,minimum width=0.58cm,minimum height=0.58cm,font=\small]
            (p\i) at (-0.66+0.66*\i,-0.365) {$\lbl$};}

    \node[ptlblg,anchor=north,font=\scriptsize] at (0,-0.66);

    \draw[ptarrow,line width=1.0pt] (0,-0.8) -- (0,-1.12);

    \node[ptghost,minimum width=2.4cm,minimum height=0.45cm,font=\scriptsize]
          (opt) at (0,-1.43) {optimizer uses $\nabla_P L$};
  \end{tikzpicture}}


  % ---------------------------------------------------------------------
% \ptwexupdate : worked-example Step 7 -- P_old moves to P_new by
% -eta * grad, shown as a short displacement in embedding space.
% Coordinate map: manifold panel -2.2..2.2 x -1.1..1.1; P_old at
% (-1.0,-0.3); P_new at (0.9,0.5); arrow between them; eta*grad label
% at the arrow midpoint.
% ---------------------------------------------------------------------
% ---------------------------------------------------------------------
% \ptwexupdate : worked-example Step 7 -- the update pipeline:
% gradient -> update rule -> new soft prompt. Vertical flow, sized to
% sit compactly beside a text column.
% Coordinate map: grad box y=1.0; update ghost box y=0.0; new-P box
% y=-1.0. Arrows fill the 0.35cm gaps between them.
% ---------------------------------------------------------------------
% ---------------------------------------------------------------------
% \ptwexupdate : worked-example Step 7 -- the update pipeline:
% gradient -> update rule -> new soft prompt. All boxes plain white,
% regular-weight text.
% Coordinate map: grad box y=1.0; update ghost box y=0.0; new-P box
% y=-1.0. Arrows fill the 0.35cm gaps between them.
% ---------------------------------------------------------------------
\newcommand{\ptwexupdate}{%
  \begin{tikzpicture}

    \uncover<2->{%

      \node[
        ptnode,
        fill=white,
        draw=ptrule,
        minimum width=2.6cm,
        minimum height=0.65cm,
        font=\small
      ]
      (grad) at (0,1.0)
      {$\eta\nabla_P L$};

      \draw[
        ptarrow,
        line width=1.0pt
      ]
      (0,0.65) -- (0,0.35);

      \node[
        ptnode,
        fill=white,
        draw=ptrule,
        minimum width=2.9cm,
        minimum height=0.65cm,
        font=\scriptsize
      ]
      (upd) at (0,0.0)
      {$P_{\text{new}} - \eta \nabla_P L$};

      \draw[
        ptarrow,
        line width=1.0pt
      ]
      (0,-0.35) -- (0,-0.65);

      \node[
        ptnode,
        fill=white,
        draw=ptrule,
        minimum width=2.9cm,
        minimum height=0.7cm,
        font=\small
      ]
      (newp) at (0,-1.0)
      {$P_{\text{new}}$};

    }

  \end{tikzpicture}%
}

  % ---------------------------------------------------------------------
% \ptworkflow : full pipeline summary -- define task, initialise P,
% freeze the model (setup row), then the forward/loss/backprop/update
% training loop, ending in the learned soft prompt.
% Coordinate map: setup row y=3.0, x=-2.8/0/2.8; loop nodes at
% (-1.4,1.6),(1.4,1.6),(1.4,0.2),(-1.4,0.2); final box y=-1.1.
% ---------------------------------------------------------------------
% ---------------------------------------------------------------------
% \ptflowfirst{current} : opening step of a multi-slide flow -- just
% the highlighted box, with a dashed arrow trailing off toward the
% next slide.
% ---------------------------------------------------------------------
\newcommand{\ptflowfirst}[1]{%
  \begin{tikzpicture}
    \node[pttuned,minimum width=4.6cm,minimum height=0.9cm,font=\small\bfseries]
          (cur) at (0,0) {#1};
    \draw[ptarrow,dashed,line width=0.9pt] (0,-0.55) -- (0,-1.05);
  \end{tikzpicture}}

% ---------------------------------------------------------------------
% \ptflowstep{previous}{current} : a middle step -- previous box faded
% at top, solid arrow down into the highlighted current box, dashed
% arrow trailing off toward the next slide.
% ---------------------------------------------------------------------
\newcommand{\ptflowstep}[2]{%
  \begin{tikzpicture}
    \node[ptghost,minimum width=4.2cm,minimum height=0.75cm,font=\small,text=ptmuted]
          (prev) at (0,1.35) {#1};
    \draw[ptarrow,line width=1.0pt] (0,0.97) -- (0,0.50);
    \node[pttuned,minimum width=4.6cm,minimum height=0.9cm,font=\small\bfseries]
          (cur) at (0,0) {#2};
    \draw[ptarrow,dashed,line width=0.9pt] (0,-0.55) -- (0,-1.05);
  \end{tikzpicture}}

% ---------------------------------------------------------------------
% \ptflowlast{previous}{current} : final step -- same as \ptflowstep
% but no trailing arrow, since the chain ends here.
% ---------------------------------------------------------------------
\newcommand{\ptflowlast}[2]{%
  \begin{tikzpicture}
    \node[ptghost,minimum width=4.2cm,minimum height=0.75cm,font=\small,text=ptmuted]
          (prev) at (0,1.35) {#1};
    \draw[ptarrow,line width=1.0pt] (0,0.97) -- (0,0.50);
    \node[pttuned,minimum width=4.6cm,minimum height=0.9cm,font=\small\bfseries]
          (cur) at (0,0) {#2};
  \end{tikzpicture}}